% Irodov Part IV §4.1 original 34--50.
% User uploaded source PDF pp.167--169, printed 158--160; answers pp.324--325.
\oscproblem{34}{158}
图4.9：质量为 $m_1=1.00$kg 的 A 物体与
$m_2=4.10$kg 的 B 物体由一轻弹簧竖直相连，
B 支承在地面上。
A 的自由竖直振动振幅为
$a=1.6$cm、角频率
$\omega=25\,\mathrm{s^{-1}}$。
求装置对地面压力的最大和最小值。
\exsolution
将 A、B 及弹簧整体视为系统，
地面支持力减总重力等于 A 物体的加速度引起的系统竖直合力：
\[
N-(m_1+m_2)g=m_1\ddot y_A.
\]
A 的加速度振幅为 $a\omega^2$，
所以
\[
\boxed{N_{\max,\min}=
(m_1+m_2)g\pm m_1a\omega^2}.
\]
代入
$(m_1+m_2)g\approx50$N、
$m_1a\omega^2=1(0.016)(625)=10$N，
\[
\boxed{N_{\max}\simeq60\,\mathrm N,\quad
N_{\min}\simeq40\,\mathrm N}.
\]
\excheck 两值均为正，说明 B 在本题所用数据下保持与地面接触。

\oscproblem{35}{158}
质量为 $m$ 的小物体放在竖直振动板上，
板自 $t=0$ 起按
$y=a(1-\cos\omega t)$ 上下运动，
$\omega=11\,\mathrm{s^{-1}}$。
求(a) 当 $a=4.0$cm 时物体压板力随时间变化；
(b) 开始脱离板的临界振幅；
(c) 物体在脱板后能上升到相对初始位置
$h=50$cm 所需的振幅。
\exsolution
在未脱离期间，板加速度
$\ddot y=a\omega^2\cos\omega t$，
竖直受力
$N-mg=m\ddot y$。
因此
\[
\boxed{N(t)=mg[1+(a\omega^2/g)\cos\omega t]}.
\]
(a) $a=0.040$m 时
$N/(mg)=1+0.494\cos(11t)$，
全程正值。
(b) 最小压力为 $m(g-a\omega^2)$，
首次可脱离条件 $a\omega^2=g$：
\[
\boxed{a_{\min}=g/\omega^2
\simeq8.1\,\mathrm{cm}}.
\]
(c) 脱板时刻满足
$\cos\omega t_*=-g/(a\omega^2)$。
那时高度
$y_*=a+g/\omega^2$，
向上速度平方
$v_*^2=a^2\omega^2-g^2/\omega^2$。
脱板后再升高度
$v_*^2/(2g)$，
总最高点
$h=a+g/(2\omega^2)+a^2\omega^2/(2g)$，
从而
\[
\boxed{a=\frac{\sqrt{2gh}-g/\omega}{\omega}
\simeq0.20\,\mathrm m}.
\]
\excheck (c) 的振幅大于(b)阈值，确保确实会离板；使用二次抛体运动处理脱离后的上升。

\oscproblem{36}{158}
将质量为 $m$ 的物体悬挂于原先未伸长的竖直弹簧上，
然后无初速度释放；弹簧劲度系数为 $k$，弹簧质量忽略。
求(a)下降位置随时间的关系；
(b)弹簧张力最大与最小值。
\exsolution
以弹簧自然长度处的释放位置为 $y=0$、
向下为正，
运动方程
$m\ddot y=mg-ky$，
令 $\omega=\sqrt{k/m}$，
初值 $y(0)=\dot y(0)=0$，
解为
\[
\boxed{y_{\rm release}(t)=\frac{mg}{k}
[1-\cos(\omega t)]}.
\]
若严格以新静平衡位置 $y_{\rm eq}=mg/k$ 为零，
则
$\boxed{y_{\rm rel.eq}(t)=-(mg/k)\cos\omega t}$。
弹簧张力 $F_s=ky_{\rm release}(t)$，
故
\[
\boxed{F_{s,\max}=2mg,\qquad
F_{s,\min}=0}.
\]
\excheck 原书文字说“相对于平衡位置的位移”，
但书末 $[1-\cos]mg/k$ 明显是相对\emph{释放位置}的下降量；
译稿同时给出两个坐标约定，避免混淆。

\oscproblem{37}{158}
质量为 $m$ 的粒子受中心回复力
$\bm F=-\alpha m\bm r$，
其中 $\alpha>0$。
初始位置 $\bm r(0)=r_0\bm i$，
速度 $\bm v(0)=v_0\bm j$。
求平面运动轨迹。
\exsolution
沿两个互相垂直方向分别满足
$\ddot x+\alpha x=0$、
$\ddot y+\alpha y=0$。
令 $\omega=\sqrt\alpha$，
使用初值：
\[
x(t)=r_0\cos\omega t,\qquad
y(t)=\frac{v_0}{\omega}\sin\omega t.
\]
消去时间：
\[
\boxed{\frac{x^2}{r_0^2}+
\frac{\alpha y^2}{v_0^2}=1}.
\]
\excheck 轨迹是以原点为中心的椭圆；
若 $v_0=\sqrt\alpha\,r_0$，特例为圆。

\oscproblem{38}{158}
质量为 $m$ 的物体悬挂在电梯顶端的轻弹簧下，
劲度系数 $k$。
从 $t=0$ 起电梯向上加速，
物体相对于电梯的初始位移及速度均为零。
求相对于原静平衡位置的运动规律：
(a) 电梯加速度 $w$ 恒定；
(b) $w(t)=\gamma t$。
\exsolution
设 $y$ 是相对电梯原静平衡位置的向下位移。
初始时重力与弹簧力平衡；
在电梯参考系，新增加的惯性力为 $mw(t)$，
所以
\[
\ddot y+\omega_0^2y=w(t),
\qquad \omega_0=\sqrt{k/m},
\quad y(0)=\dot y(0)=0.
\]
(a) 对常数 $w$：
\[
\boxed{y(t)=\frac w{\omega_0^2}
(1-\cos\omega_0t)}.
\]
(b) 用受迫系统格林函数
$y(t)=\omega_0^{-1}\int_0^t
w(s)\sin[\omega_0(t-s)]\,\mathrm ds$，
当 $w=\gamma t$ 时
\[
\boxed{y(t)=\frac{\gamma}{\omega_0^3}
(\omega_0t-\sin\omega_0t)}.
\]
\excheck 两式初始位置、速度为零，
代回微分方程均得到相应驱动加速度。

\oscproblem{39}{158}
质量 $m=0.50$kg 的物体悬挂于弹性系数
$k=50\,\mathrm{N/m}$ 的橡皮绳下。
向下最多拉离平衡位置多少，释放后振动才仍是谐振动？
此临界振动能量是多少？
\exsolution
橡皮绳只在伸长时产生张力，
新静平衡伸长量
$x_0=mg/k$。
若围绕此点振幅 $A>x_0$，
上转折点绳将松弛，回复力不再线性。
临界振幅
\[
\boxed{A_{\max}=\frac{mg}{k}
=\frac{0.50(9.8)}{50}
=9.8\,\mathrm{cm}\approx10\,\mathrm{cm}}.
\]
此时振动能量
\[
\boxed{E_{\max}=\frac12kA_{\max}^2
=\frac{m^2g^2}{2k}\simeq0.24\,\mathrm J}.
\]
\excheck 本版书后解析式与上式一致，
但末尾数值写成约 $4.8\,\mathrm{mJ}$，
与给定 $m,k$ 不符，差约50倍；记录待核验。

\oscproblem{40}{158--159}
图4.10：质量为 $m$ 的小球从静止自由落下高度 $h$，
撞在弹簧秤的无质量托盘上并粘住，
弹簧刚度为 $k$，初始未压缩。
求其后的竖直振动振幅和机械振动能量。
\exsolution
接触前速度
$v_0=\sqrt{2gh}$。
接触后新静平衡弹簧压缩量
$x_0=mg/k$，
相对于新平衡位置初始坐标
$y(0)=-x_0$，初速 $v_0$ 向下。
角频率 $\omega_0=\sqrt{k/m}$，
振幅由初始相空间半径给
\[
A^2=y(0)^2+\frac{v_0^2}{\omega_0^2}
=\left(\frac{mg}{k}\right)^2+\frac{2mgh}{k}.
\]
所以
\[
\boxed{A=\frac{mg}{k}
\sqrt{1+\frac{2hk}{mg}}},
\quad
\boxed{E=\frac12kA^2
=mgh+\frac{m^2g^2}{2k}}.
\]
\excheck 题设托盘无质量，粘住碰撞后球速度仍为
$\sqrt{2gh}$；若托盘有质量则碰撞时损失动能。

\oscproblem{41}{159}
承接第40题，现托盘质量为 $M$，
在物体落下前托盘已停在自身的静平衡位置。
小球质量 $m$ 从高度 $h$ 落下后粘住托盘。
求碰撞后共同振动的振幅。
\exsolution
小球落下速度 $v_0=\sqrt{2gh}$。
完全非弹性粘合碰撞，沿竖直方向短时动量守恒，
新组合体初速
$v_+=m v_0/(m+M)$。
托盘原来压缩 $Mg/k$，
粘合后新平衡压缩
$(M+m)g/k$，
因此相对新平衡初始坐标
$y_0=-mg/k$。
组合质量为 $m+M$，角频率
$\omega_0=\sqrt{k/(m+M)}$，
故
\[
\boxed{A=
\sqrt{\left(\frac{mg}{k}\right)^2+
\frac{m^2(2gh)}{k(m+M)}}
=\frac{mg}{k}
\sqrt{1+\frac{2hk}{(m+M)g}}}.
\]
\excheck $M\to0$ 时精确恢复第40题；碰撞中机械能一般不守恒，必须先守恒动量。

\oscproblem{42}{159}
质量为 $m$ 的粒子在 $xy$ 平面内运动。
原书题面印作速度相关力
$\bm F=\alpha(\dot y\bm i+\dot x\bm j)$，
$\alpha>0$，
初始 $\bm r(0)=0$、$\bm v(0)=v_0\bm j$。
求 $x(t),y(t)$ 及轨迹。
\exsolution
按\emph{原版题面两个正号}，令
$\omega=\alpha/m$，
运动方程为
$\ddot x=\omega\dot y$、
$\ddot y=\omega\dot x$。
初值 $\dot x(0)=0$、$\dot y(0)=v_0$，
解为
$\dot x=v_0\sinh\omega t$、
$\dot y=v_0\cosh\omega t$，
故
\[
\boxed{x(t)=\frac{v_0}{\omega}
(\cosh\omega t-1),\qquad
y(t)=\frac{v_0}{\omega}\sinh\omega t}.
\]
轨迹为
\[
\boxed{\left(1+\frac{\omega x}{v_0}\right)^2-
\left(\frac{\omega y}{v_0}\right)^2=1}.
\]
\emph{书末答案采用另一种力：}
它写
$\ddot y=-\omega\dot x$，
这对应
$\bm F=\alpha(\dot y\bm i-\dot x\bm j)$，
此时才是
$x=(v_0/\omega)(1-\cos\omega t)$、
$y=(v_0/\omega)\sin\omega t$
的圆形轨迹。
\excheck 不能把原题印出的“$+$”悄然改成“$-$”；
本稿保留字面题面的双曲线解及书末另一个符号模型的圆轨迹，待校对印刷源。

\oscproblem{43}{159}
图4.11：轻质薄壁球壳半径 $R$，
内部充满水，球心以轻刚杆悬于点 O，距离为 $\ell$。
水结冰后小振动周期改变多少倍？
忽略水的黏度与结冰体积变化。
\exsolution
结冰前，水在球壳内近似平动而不随壳整体刚性自转，
系统可视为质量为 $m$、有效摆长为 $\ell$ 的单摆，
$J_{\rm liquid}=m\ell^2$。
结冰后水与壳刚性转动，实心冰球绕球心转动惯量
$J_c=(2/5)mR^2$，
绕悬点转动惯量
$J_{\rm ice}=m\ell^2+(2/5)mR^2$。
两态重力回复力矩系数相同为 $mg\ell$，
小摆周期 $T=2\pi\sqrt{J/(mg\ell)}$。
故
\[
\boxed{\frac{T_{\rm ice}}{T_{\rm liquid}}
=\sqrt{1+\frac25\left(\frac R\ell\right)^2}}.
\]
\excheck 比值大于1，结冰后因新增自转惯量而振动变慢。

\oscproblem{44}{159}
图4.12：一根质量 $m$、长 $\ell$ 的均匀细杆，
上端 O 铰接，自然垂直向下；
下端由合刚度为 $k$ 的水平弹簧保持。
求杆的微小平面摆动角频率。
\exsolution
杆相对于 O 转动惯量
$J=m\ell^2/3$。
小转角 $\theta$ 时，重力回复力矩
$\tau_g\simeq-(mg\ell/2)\theta$；
杆底横移 $\ell\theta$，
弹簧回复力 $-k\ell\theta$，
相对 O 的回复力矩
$\tau_s\simeq-k\ell^2\theta$。
所以
\[
J\ddot\theta+
\left(\frac{mg\ell}{2}+k\ell^2\right)\theta=0,
\]
\[
\boxed{\omega_0=
\sqrt{\frac{3g}{2\ell}
\left(1+\frac{2k\ell}{mg}\right)}}.
\]
\excheck $k\to0$ 时还原铰接细杆的物理单摆角频率 $\sqrt{3g/(2\ell)}$。

\oscproblem{45}{159}
图4.13：质量 $m=1.5$kg 的均匀水平杆由两条
长 $\ell=90$cm 的相同细绳悬吊，
绕竖直中心轴转动，
两绳偏离竖直的最大角度
$\alpha=5.0^\circ$，然后释放。
求(a)小转角振动周期；(b)振动能量。
\exsolution
设杆长（两悬点间距）为 $d$，
绕竖直中心轴转动惯量
$J=md^2/12$。
小转角 $\varphi$ 使两绳的横向偏移量约为 $d\varphi/2$，
对应绳倾角 $\alpha\simeq d\varphi/(2\ell)$，
支承高度每条绳缩短
$\Delta h\simeq\ell\alpha^2/2$，
重力势能增加
$U\simeq mg\ell\alpha^2/2
=mgd^2\varphi^2/(8\ell)$。
故扭转回复刚度 $k_\varphi=mgd^2/(4\ell)$，
\[
\boxed{T=2\pi\sqrt{\frac{J}{k_\varphi}}
=2\pi\sqrt{\frac{\ell}{3g}}
\simeq1.1\,\mathrm s}.
\]
最大偏角 $\alpha$ 时全部能量为
\[
\boxed{E=\frac12mg\ell\alpha^2
\simeq0.050\,\mathrm J}.
\]
\excheck $d$ 在周期中消去；$\alpha$ 计算时应换成弧度。

\oscproblem{46}{160}
图4.14：均匀水平圆盘质量为 $m$、半径 $R$，
用扭转系数为 $k$ 的细杆悬挂。
初始扭转角为 $\varphi_0$，
同时给予有向角速度 $\dot\varphi_0$。
求扭转小振动振幅及能量。
\exsolution
圆盘相对中心竖轴的转动惯量
$J=mR^2/2$，
运动方程
$J\ddot\varphi+k\varphi=0$，
频率
$\omega_0=\sqrt{k/J}$。
由初始相空间关系
$\varphi_m^2=\varphi_0^2+
\dot\varphi_0^2/\omega_0^2$，
得到
\[
\boxed{\varphi_m=
\sqrt{\varphi_0^2+\frac{mR^2}{2k}
\dot\varphi_0^{\,2}}},
\quad
\boxed{E=\frac12k\varphi_m^2}.
\]
\excheck 原书对初始角与角速度的文字排印符号接近；
译稿用 $\varphi_0$ 与 $\dot\varphi_0$ 明确区分。

\oscproblem{47}{160}
均匀细杆质量 $m$、长度 $\ell$，
绕上端水平轴作小角振动。
初始偏角为 $\theta_0$、
角速度为 $\dot\theta_0$。
求平均一个周期内的动能。
\exsolution
杆转动惯量
$J=m\ell^2/3$，
重力小角回复系数
$K=mg\ell/2$。
振动总能量
$E=\tfrac12J\dot\theta_0^2+
\tfrac12K\theta_0^2$。
简谐振动中一周期平均动能等于总能量的一半，
所以
\[
\boxed{\langle T_k\rangle=
\frac14J\dot\theta_0^2+
\frac14K\theta_0^2
=\frac{m\ell^2\dot\theta_0^2}{12}
+\frac{mg\ell\theta_0^2}{8}}.
\]
\excheck 小角近似中 $\theta_0$ 须以弧度计；$E/2$ 的等分仅指完整周期平均。

\oscproblem{48}{160}
物理单摆的质心高于悬挂点，
起初处于倒立不稳定平衡附近。
它开始倒下、经过最低稳定平衡位置时的角速度为 $\Omega$。
忽略摩擦，求最低位置附近小角振动周期。
\exsolution
设悬点到质心距离为 $h$，
转动惯量为 $J$。
从上方倒立平衡到下方稳定平衡，重力势能下降
$2mgh$，所以
$\tfrac12J\Omega^2=2mgh$，
即
$mg h/J=\Omega^2/4$。
下方小振动角频率
$\omega_0=\sqrt{mgh/J}=\Omega/2$，
所以
\[
\boxed{T=\frac{2\pi}{\omega_0}
=\frac{4\pi}{\Omega}}.
\]
\excheck 严格在不稳定平衡点零初速会永远不动；
按题意需有任意微小扰动使其开始下落。

\oscproblem{49}{160}
物理单摆原来的小振动角频率
$\omega_1=15.0\,\mathrm{s^{-1}}$；
在距悬点下方 $\ell=20$cm 处固定一小物体，
质量 $m=50$g，
新角频率变成
$\omega_2=10.0\,\mathrm{s^{-1}}$。
求原单摆对悬点的转动惯量。
\exsolution
原摆的重力回复力矩系数设为 $K$，
有
$K=J\omega_1^2$。
增加小质量后
$K'=K+mg\ell$，
$J'=J+m\ell^2$。
新频率满足
\[
\omega_2^2=
\frac{J\omega_1^2+mg\ell}{J+m\ell^2}.
\]
解出
\[
\boxed{J=
\frac{m\ell^2(\omega_2^2-g/\ell)}
{\omega_1^2-\omega_2^2}}
\simeq8.2\times10^{-4}\,\mathrm{kg\,m^2}
=0.82\,\mathrm{g\,m^2}}.
\]
\excheck 因 $J>0$，需
$\omega_2^2>g/\ell$；本题数值满足。

\oscproblem{50}{160}
两个物理单摆绕同一水平轴作小振动，
角频率分别为 $\omega_1,\omega_2$，
转动惯量为 $J_1,J_2$。
在稳定平衡位置把它们刚性连在一起，
求复合单摆频率。
\exsolution
单摆小角回复力矩系数
$K_i=J_i\omega_i^2$。
刚性连接后两摆偏角相同，
总转动惯量
$J=J_1+J_2$，
总回复系数
$K=K_1+K_2$。
所以
\[
\boxed{\omega=
\sqrt{\frac{J_1\omega_1^2+J_2\omega_2^2}
{J_1+J_2}}}.
\]
\excheck $\omega^2$ 是两频率平方的惯量加权平均，
因此 $\min(\omega_1,\omega_2)\le\omega\le\max(\omega_1,\omega_2)$。
